\scanpage{12}
% 1–2
Les conditions a) b) c) sont les \uline{conditions algébriques} élémentaires. On veut de plus

% 3, left
% slanted
\margin{après d) (condition différentielle)}

% 3–10
\deleted{\unsure{d}}\unsure{e)|d)} (Condition hensélienne) Si $u:X_n-X_n$ est un $A_*$-morphisme qui est un isomorphisme formel (i.e. si $f_1,\ldots,f_n\in A_n$ sont tels que $\varepsilon(f_1)\ldots\varepsilon(f_n)$ sont nilpotents dans $k$ et
\[
 \varepsilon\left(\det\left(\frac{\partial f_i}{\partial t_j}\right)\right)
 =\text{dét}\left(\varepsilon\left(\frac{\partial f_i}{\partial x_j}\right)\right)
 =\det\left(\frac{\partial f_i}{\partial\unsure{x_j|j|t_j}}(0)\right)
\]
est un élément inv. de $k$) alors $u^{-1}$ est aussi un $A_*$-morphisme ($\Leftrightarrow$ le foncteur des \deleted{$A_*$}schémas $A_*$ formels vers les schémas formels sur $k$ est conservatif).

% 11–16
\uline{Cor. 1}\qquad Si $f\in A_n$ et $f$ inversible dans $S_n$, i.e. $\varepsilon(f)=f(0)$ inversible dans $k$, alors $f^{-1}\in A_n$ ; \inserted{$\Longleftrightarrow$} \deleted{En particulier \inserted{\deleted{ou encore}}} si $f\in A_n$, $\varepsilon(f)=0$, alors $(1+f)^{-1}\in A_n$ [$\Longleftrightarrow$ si $f\in A_n$ et $\varepsilon(f)=1$, alors $f^{-1}\in A_n$ $\}\Longleftrightarrow A_n\cap\mathfrak m_n\subset\mathrm{rad}\,A_n$, où \deleted{$\mathfrak m_n=$} $\mathfrak m_n=(t_1,\ldots,t_n)S_n$ $\}$ \deleted{\unsure{au radical de $A_n\cap\mathfrak m_n$}}]

% 13, left
% slanted
\margin{moyennant condition d) (différentielle)}

% 17–20
En effet, considérons la \deleted{\unsure{multi}} l'endomorphisme
\[
 u:(x_1,\ldots,x_n,x_{n+1})\mapsto(x_1,\ldots,x_n,f(x_1,\ldots,x_n).x_{n+1})
\]
de $X_{n+1}$, il est un isom. formel d'inverse
\[
 (x_1,\ldots,x_n,x_{n+1})\mapsto(x_1,\ldots,x_n,f(x_1,\ldots,x_n)^{-1}x_{n+1}),
\]

% 21–24
donc comme \deleted{c'est}\inserted{$u$ est} un $A$-morphisme, par d) il en est de $=$ de $u^{-1}$, donc $g(t_1,\ldots,t_{n+1})=f(t_1,\ldots,t_n)^{\unsure{-1|-2}}t_{n+1}\in A_{n+1}$, \deleted{i.e. $X_{n+1}\to X_1$ \unsure{est un}} donc $g(t_1,\ldots,t_n,1)\in A_n$ i.e. $f^{-1}\in A_n$, cqfd.

% 22, left
% slanted
\margin{résulte de la condition différentielle}

% 24, left
% slanted
\margin{\unsure{incorrect|inconvénient}, $1$ n'est pas nilpotent ! Il faut \unsure{obliger|exiger} une propriété de divisibilité $t_{n+1}f\in A_{n+1}\Longrightarrow f\in A_{n+1}$}

% 25–27
\uline{Cor. 2}\qquad Soit $\inserted{P(Z)=}Z^N+f_1 Z^{N-1}+\ldots+f_N\in A_n[Z]$ polynôme unitaire, soit $z\in k_0$ tel que
\[
 P^\varepsilon(z)=z^N+\deleted{\unsure{f}}a_1 z^{N-1}\ldots+a_N=0,\quad
 \text{où }\deleted{\unsure{\xi}}\ a_i=\varepsilon(f_i)=f_i(0)\in k_0,
\]
% archive sheet number: 11
